\chapter{Assegnamento}

Sia dato un grafo bipartito completo $G = (V_1 \cup V_2, E)$,
nel problema di assegnamento si cerca di effettare un matching fra i due sotto insieme di 
vertici di modo da ottenere il costo minore possbili per ogni accoppiamento.


\[
\begin{array}{c|cccc}
      & \text{Job 1} & \text{Job 2} & \text{Job 3} & \text{Job 4} \\
\hline
\text{Worker 1} & 10 & 8 & 9 & 7 \\
\text{Worker 2} & 6 & 9 & 7 & 5 \\
\text{Worker 3} & 8 & 7 & 6 & 9 \\
\text{Worker 4} & 5 & 6 & 7 & 8 \\
\end{array}
\]


\section{Algoritmo ungherese}
\subsection{Riduzione}
Il primo passo dell'algoritmo ungherese \`e quello di ridurre al minimo la matrice dei costi e si porsegue nella seguente maniera:

\begin{enumerate}
  \item Per ogni riga, sottrarre il minimo valore della riga a tutti i valori della riga.
  \item Per ogni colonna, sottrarre il minimo valore della colonna a tutti i valori della colonna.
\end{enumerate}


Applicando la riduzione alla matrice dei costi si ottiene:

Sottrai il valore minimo di ogni riga dalla riga corrispondente:
\begin{enumerate}
   \item Sottrai 7 da ogni valore nella riga del Worker 1
   \item Sottrai 5 da ogni valore nella riga del Worker 2
   \item Sottrai 6 da ogni valore nella riga del Worker 3
   \item Sottrai 5 da ogni valore nella riga del Worker 4
\end{enumerate}

\[
\begin{array}{c|cccc}
      & \text{Job 1} & \text{Job 2} & \text{Job 3} & \text{Job 4} \\
\hline
\text{Worker 1} & 3 & 1 & 2 & 0 \\
\text{Worker 2} & 1 & 4 & 2 & 0 \\
\text{Worker 3} & 2 & 1 & 0 & 3 \\
\text{Worker 4} & 0 & 1 & 2 & 3 \\
\end{array}
\]

\begin{enumerate}
   \item Sottrai 0 da ogni valore nella colonna del Job 1
   \item Sottrai 1 da ogni valore nella colonna del Job 2
   \item Sottrai 0 da ogni valore nella colonna del Job 3
   \item Sottrai 0 da ogni valore nella colonna del Job 4
\end{enumerate}

\[
\begin{array}{c|cccc}
      & \text{Job 1} & \text{Job 2} & \text{Job 3} & \text{Job 4} \\
\hline
\text{Worker 1} & 3 & 0 & 2 & 0 \\
\text{Worker 2} & 1 & 3 & 2 & 0 \\
\text{Worker 3} & 2 & 0 & 0 & 3 \\
\text{Worker 4} & 0 & 0 & 2 & 3 \\
\end{array}
\]

\subsection{Fase di assegnazione}
Dopo aver ridotto la matrice, cerco di coprire con varie righe verticali ed orizzontali la matrice, le righe vanno messe dove sono gli zeri all'interno della matrice ridotta.


Se riesco a coprire tutti gli zeri con $n$ linee, dove $n$ rappresenta il numero di righe e colonne della matrice ridotta, possono selezionare i vari elementi dove sono preseti gli zeri e questi rappresentano l'assegnamento ottimo.
Al che seleziono i valori della matrice originale in corrispondenza degli zeri selezionati nella matrice ridotta.


Se non dovessi avere abbastanza zeri da poter coprire tutta la matrice, allora devo procedere con la fase di correzione.

Si seleziona l'elemento pi\`u piccolo diverso da zero della riga o colonna incriminata e si sottrae a tutti gli elementi della riga o colonna, mentre si aggiunge il valore selezionato nella cella dove \`e presente un incrocio di riga e colonna.



